\documentclass[11pt]{article}
\usepackage[margin=1in]{geometry}
\usepackage{amsmath,amssymb,amsthm,mathtools}
\usepackage[T1]{fontenc}
\usepackage{hyperref,booktabs,enumitem}
\hypersetup{colorlinks=true,linkcolor=blue,citecolor=blue,urlcolor=blue}
\title{Detector-Gated Bounds for Multi-Band Dirichlet Polynomials}
\author{Shawn Calvin Snelling\\\small AXEZENT AI Research}
\date{Research manuscript --- version 0.2.0\\October 23, 2026}
\newtheorem{theorem}{Theorem}[section]
\newtheorem{lemma}[theorem]{Lemma}
\newtheorem{proposition}[theorem]{Proposition}
\theoremstyle{remark}\newtheorem{remark}[theorem]{Remark}
\begin{document}
\maketitle
\begin{abstract}
We establish a finite-grid witness principle for finite Dirichlet polynomials and record an exact finite Fourier-energy identity for Möbius exponential sums. The first result reduces a continuous-time threshold question on a compact interval to finitely many evaluations, provided an explicit derivative bound and grid covering radius are supplied. The second rewrites a finite Fourier energy as a sinc-weighted sum of arithmetic correlations. We state exactly what these identities establish and what they do not: neither identity alone supplies a uniform asymptotic estimate, a limit-exchange theorem, or a proof of the Riemann Hypothesis.
\end{abstract}

\section{Scope and notation}
A finite Dirichlet polynomial is
\[
P(t)=\sum_{n\in B}a_n n^{-it},\qquad n^{-it}:=e^{-it\log n},
\]
where $B$ is a finite set of positive integers and $a_n\in\mathbb C$. All interval statements concern a fixed compact interval $I=[T_0,T_1]$. A detector is simply a rule that tests values against a specified threshold; it adds no unproved mathematical guarantee.

\section{Derivative and Lipschitz bounds}
\begin{lemma}[Derivative bound]\label{lem:deriv}
For $P(t)=\sum_{n\in B}a_n n^{-it}$, put
\[
L:=\sum_{n\in B}|a_n|\log n.
\]
Then $|P'(t)|\le L$ for every real $t$.
\end{lemma}
\begin{proof}
The sum is finite, so termwise differentiation gives
$P'(t)=-i\sum_{n\in B}a_n(\log n)n^{-it}$.
Since $|n^{-it}|=1$, the triangle inequality yields $|P'(t)|\le L$.
\end{proof}
\begin{corollary}\label{cor:lip}
For all real $t,u$, $|P(t)-P(u)|\le L|t-u|$.
\end{corollary}
\begin{proof}
Integrate $P'$ along the segment between $u$ and $t$ and apply Lemma~\ref{lem:deriv}.
\end{proof}

\section{Finite-grid witness theorem}
\begin{definition}
A finite set $G\subset I$ has covering radius at most $h$ on $I$ if every $t\in I$ lies within distance $h$ of some $g\in G$.
\end{definition}
\begin{theorem}[Finite-grid witness lemma]\label{thm:grid}
Let $P$, $I$, and $G$ be as above, with covering radius at most $h$. If $\tau>0$, $0<\varepsilon<\tau$, and $Lh\le\varepsilon$, then
\[
\sup_{t\in I}|P(t)|\ge\tau
\quad\Longrightarrow\quad
\max_{g\in G}|P(g)|\ge\tau-\varepsilon.
\]
Equivalently, if $|P(g)|<\tau-\varepsilon$ for every $g\in G$, then $|P(t)|<\tau$ for every $t\in I$.
\end{theorem}
\begin{proof}
Continuity and compactness give a point $t_*\in I$ at which $|P|$ attains its maximum. Choose $g_*\in G$ with $|t_*-g_*|\le h$. By the reverse triangle inequality and Corollary~\ref{cor:lip},
\[
|P(g_*)|\ge |P(t_*)|-|P(t_*)-P(g_*)|
\ge\tau-Lh\ge\tau-\varepsilon.
\]
For the contrapositive, for each $t\in I$ choose $g\in G$ with $|t-g|\le h$. Then $|P(t)|\le |P(g)|+Lh<\tau-\varepsilon+\varepsilon=\tau$.
\end{proof}
\begin{remark}[Spacing convention]
If a grid includes both endpoints and its largest adjacent spacing is $\Delta$, its covering radius is at most $\Delta/2$. Thus $\Delta\le2\varepsilon/L$ suffices when $L>0$. For nonuniform or one-sided grids, verify the actual covering radius; do not silently identify it with the largest adjacent spacing. If $L=0$, $P$ is constant.
\end{remark}

\section{Numerical certificates}
If a numerical routine returns $\widetilde P(g)$ and a rigorous error bound $|P(g)-\widetilde P(g)|\le\delta$ holds at every grid point, then
\[
\max_{g\in G}\bigl(|\widetilde P(g)|+\delta\bigr)+Lh<\tau
\]
is a sufficient certificate that $|P(t)|<\tau$ throughout $I$. Without a rigorous evaluation-error bound, a floating-point scan is numerical evidence, not a proof certificate.

A reproducible certificate should record: (i) the index set and coefficient encoding; (ii) interval endpoints, threshold, tolerance, and $L$; (iii) the grid-generation rule and verified $h$; (iv) arithmetic, precision, rounding policy, and error bound; (v) sampled values or a reproducible regeneration procedure; (vi) the final checked inequality; and (vii) source revision, compiler/runtime, and hashes. C++ workers may generate candidate data, but correctness follows from verifying the stated hypotheses and inequalities, not from speed alone.

\section{A separate exact Möbius Fourier-energy identity}
For $N\ge1$, define
\[
F_N(u):=\frac1{\sqrt N}\sum_{n\le N}\mu(n)e^{2\pi inu/N},\qquad
Q_N(c):=\frac1{2c}\int_{-c}^c|F_N(u)|^2\,du
\]
for $c>0$. For $0\le d<N$, let
\[
C_N(d):=\sum_{\substack{n\ge1\\n+d\le N}}\mu(n+d)\mu(n),
\]
and set $C_N(-d)=C_N(d)$. Let $\operatorname{sinc}(x)=\sin(x)/x$ for $x\ne0$ and $\operatorname{sinc}(0)=1$.
\begin{proposition}[Exact sinc-correlation identity]\label{prop:sinc}
For all $N\ge1$ and $c>0$,
\[
Q_N(c)=\frac1N\sum_{m,n\le N}\mu(m)\mu(n)
\operatorname{sinc}\!\left(\frac{2\pi c(m-n)}N\right)
=\frac1N\sum_{|d|<N}C_N(d)
\operatorname{sinc}\!\left(\frac{2\pi cd}N\right).
\]
Consequently,
\[
Q_N(c)=\frac1N\sum_{n\le N}\mu(n)^2+
\frac2N\sum_{1\le d<N}C_N(d)
\operatorname{sinc}\!\left(\frac{2\pi cd}N\right).
\]
\end{proposition}
\begin{proof}
Expand the finite product:
\[
|F_N(u)|^2=\frac1N\sum_{m,n\le N}\mu(m)\mu(n)e^{2\pi i(m-n)u/N}.
\]
The sum is finite, so termwise integration is valid. Direct integration gives
\[
\frac1{2c}\int_{-c}^c e^{2\pi i d u/N}\,du
=\operatorname{sinc}\!\left(\frac{2\pi cd}N\right).
\]
This proves the double-sum formula. Grouping terms by $d=m-n$ gives the correlation formula. The diagonal is the $d=0$ term, and the terms for $d$ and $-d$ agree, proving the final expression.
\end{proof}
\begin{remark}
This is an exact finite identity. It does not imply a bound uniform in $N$, nor does it justify exchanging a limit with a sum or integral. Any asymptotic assertion such as $Q_N(c)=N^{o(1)}$ requires additional estimates not proved by this identity.
\end{remark}

\section{Connecting the tools without conflating them}
The Dirichlet polynomial $P(t)$ uses phases $e^{-it\log n}$, whereas $F_N(u)$ uses phases $e^{2\pi inu/N}$. These are different functions. The grid theorem may be applied to $F_N$ as its own finite exponential polynomial, after deriving a derivative bound for that function.

In fact,
\[
F_N'(u)=\frac{2\pi i}{N\sqrt N}\sum_{n\le N}n\mu(n)e^{2\pi inu/N},
\]
and hence
\[
|F_N'(u)|\le\frac{2\pi}{N\sqrt N}\sum_{n\le N}n|\mu(n)|
\le\frac{\pi(N+1)}{\sqrt N}.
\]
This gives a valid, though potentially crude, finite-interval grid bound for $F_N$. It gives no cancellation estimate for the Möbius coefficients and no asymptotic control as $N\to\infty.

A complete argument must keep four steps distinct: (1) exact finite algebra; (2) a finite-interval grid certificate with derivative and evaluation errors checked; (3) any required uniform asymptotic correlation or tail estimate; and (4) the separate limiting implication that uses that estimate. Steps (3) and (4) are not supplied here.

\section{Proof-status ledger}
\begin{center}
\begin{tabular}{p{0.16\linewidth}p{0.57\linewidth}p{0.17\linewidth}}
\toprule
ID & Statement & Status\\
\midrule
DG-001 & Derivative bound for finite Dirichlet polynomials & Proved\\
DG-002 & Lipschitz estimate & Proved\\
DG-003 & Finite-grid witness lemma & Proved\\
MF-001 & Finite Möbius Fourier-energy identity & Proved\\
ASY-001 & Uniform asymptotic cancellation of sinc-weighted correlations & Open here\\
LIM-001 & Any required limit-exchange or uniform-tail step & Not established\\
RH-001 & Deduction of RH from the statements in this paper & Not established\\
\bottomrule
\end{tabular}
\end{center}

\section{Conclusion}
The finite-grid witness lemma rigorously converts a continuous-time threshold question on a compact interval into finite evaluations when all hypotheses and numerical errors are controlled. The Möbius Fourier-energy identity exposes an exact sinc-weighted correlation sum. Connecting these results to an asymptotic theorem requires additional uniform estimates and a separately justified limiting argument. No global asymptotic closure or proof of the Riemann Hypothesis is claimed here.

\end{document}
